The modern electricity grid operates through wholesale markets in which
producers and large consumers interact---directly or through a central
operator---to determine how much power is generated and consumed.  In a
stylised day-ahead market, producers submit offers specifying the quantities of
electricity they are willing to supply at various price points, while suppliers
(and possibly a handful of very large consumers) submit corresponding demand
bids.  A system operator then clears the system by finding the price at which
aggregate supply equals aggregate demand, setting a single clearing price for
each settlement period (typically one hour).  This mechanism promotes economic
efficiency by rewarding low-cost producers and incentivising demand
response~\parencite{stoft2002power, kirschen2018fundamentals}.

Real markets differ considerably in how centralised this process is.  The
Israeli market, which is the setting of this study, is comparatively
centralised: the system operator (Noga) forecasts national demand for each hour
of the following day, and this day-ahead forecast---together with the resulting
price signal---determines how much electricity producers are scheduled to
deliver in each settlement period.  The process is therefore less the
continuously cleared double auction of textbook descriptions and more a
coordinated, forecast-driven dispatch.  Regardless of the precise market design,
however, the day-ahead demand forecast is the pivotal quantity: it is fixed
before real-time demand is known, and its accuracy propagates directly into
generation scheduling and price.

A central challenge in this architecture is that procurement decisions must be
finalised a day in advance, before actual real-time demand is known.  The
accuracy of demand forecasts therefore directly affects the operational
decisions and financial outcomes of all market participants.  It remains an open
question to what extent individual suppliers maintain their own demand models.  A
plausible account of current practice is that, once the operator publishes its
day-ahead demand forecast, each supplier simply draws its own share from that
forecast based on experience, and the market settles on a relatively stable price
that is reasonable on average but may embed sizeable, systematic error margins.
If so, more sophisticated forecasting---whether by the suppliers or by the system
operator---could translate into lower retail prices, which is part of what
motivates the present work.

A system operator that accurately anticipates demand can procure the right
amount of energy at competitive prices.  Inaccurate forecasts, by contrast,
force participants into costly corrective actions in the balancing or real-time
markets, where prices can be dramatically different from day-ahead levels.  It is
worth stressing that these price differences arise from two distinct sources of
uncertainty: errors in predicting \emph{demand}, and errors in predicting
available \emph{production capacity}---the latter being especially pronounced in
systems that rely heavily on variable renewable generation such as solar and
wind~\parencite{morales2014integrating}.  This paper is concerned exclusively
with the first source, uncertainty in day-ahead demand; production-side
uncertainty, while important, is outside our scope.

\subsection{Asymmetric Costs in Electricity Markets}

A critical and often overlooked aspect of electricity demand forecasting is the
\emph{asymmetry} of the costs associated with forecast errors.  When demand is
\emph{overestimated}, the system operator or retailer procures more electricity
than is ultimately needed.  The excess supply can typically be sold back into the
balancing market, curtailed, or stored (if storage is available), resulting in a
relatively modest financial loss.  In deregulated markets with well-functioning
balancing mechanisms, the per-unit cost of over-procurement is generally bounded
and predictable.

When demand is \emph{underestimated}, the cost of error is usually much higher.
The system operator must source additional electricity in the real-time market at
short notice to prevent blackouts.  Real-time prices are typically highly
volatile and can spike to several multiples of the day-ahead price during periods
of scarcity.  In extreme cases, the system may lack sufficient reserve capacity
to meet demand, triggering involuntary load shedding, emergency imports, or the
activation of expensive peaking units.  Beyond the direct financial costs, supply
shortfalls carry significant reliability and regulatory
implications~\parencite{stoft2002power, kirschen2018fundamentals}.

Empirical studies in various electricity markets suggest that the ratio of
under-procurement costs to over-procurement costs, $c_u / c_o$, typically lies in
the range $[5,\,20]$, depending on market structure, regulatory context, and the
availability of real-time resources~\parencite{zhang2022costoriented}.  This
asymmetry fundamentally changes the optimal forecasting problem: rather than
seeking a prediction that minimises expected squared error, a rational market
participant should seek a prediction that minimises expected asymmetric cost.

This decision-theoretic view has developed alongside a substantial body of work
on electricity forecasting more broadly.  \textcite{weron2014electricity}
provides a comprehensive review of the state of the art in electricity price
forecasting, surveying statistical, machine-learning, and hybrid approaches.
While its focus is on price rather than demand, the review highlights the central
role of model selection, loss-function choice, and benchmark comparisons that
also apply in the load-forecasting context.

\textcite{hong2016probabilistic} offer a tutorial review of probabilistic
electric load forecasting, covering methods ranging from regression-based
quantile estimation to ensemble and Bayesian approaches.  Crucially, they argue
that point forecasts are insufficient for operational decision-making and
advocate for full predictive distributions or at least prediction intervals---a
view consistent with the quantile-forecasting objective we adopt.

The practical importance of probabilistic evaluation metrics, and of the pinball
loss in particular, was cemented by the Global Energy Forecasting Competition
2014~\parencite{hong2016gefcom2014}.  GEFCom2014 was the first major forecasting
competition to use pinball loss as its primary scoring rule across the
electricity load, price, wind, and solar tracks.  It established pinball loss as a
community standard and demonstrated that models explicitly optimised for quantile
accuracy substantially outperform those adjusted only post-hoc.

Despite the theoretical clarity of this framing, the majority of practical
forecasting systems in electricity markets continue to use symmetric loss
functions---most commonly MSE or mean absolute error (MAE)---both for training
predictive models and for evaluating their performance.  This disconnect between
training objective and operational cost constitutes a systematic source of
suboptimality that has received growing attention in the academic
literature~\parencite{wang2019pinball, lin2020asymmetric, wu2021svr,
zhang2022costoriented} but has yet to see widespread adoption in practice.

\subsection{Cost-Oriented and Asymmetric Learning for Load Forecasting}

The theoretical foundation for prediction under asymmetric costs was laid by
\textcite{koenker1978regression}, who introduced \emph{regression quantiles} as a
natural generalisation of sample quantiles to the linear model.  They showed that
minimising the asymmetric pinball loss (equivalently, the ``check function'')
yields an estimator of the conditional $\tau$-quantile of the response, and
established its asymptotic properties.  This result directly underpins our use of
a cost-sensitive asymmetric loss function---instantiated as a non-linear variant
of the pinball loss---as the training objective for day-ahead demand forecasting.

\textcite{christoffersen1997optimal} provided the corresponding
decision-theoretic foundation.  They showed that under a general asymmetric loss
function the optimal point forecast is \emph{biased}, and that the optimal amount
of bias is time-varying and depends on higher-order conditional moments such as
conditional variance.  In particular, their analysis implies that any forecast
trained with a symmetric objective (e.g.\ MSE) will be systematically suboptimal
when the true cost structure is asymmetric, even if the conditional mean is
estimated perfectly.

Building on these foundations, a more recent line of work addresses the mismatch
between symmetric training objectives and asymmetric operational costs directly
in the electricity setting, embedding the cost asymmetry into the learning
objective itself rather than correcting for it after the fact.

\textcite{wang2019pinball} train a Long Short-Term Memory (LSTM) network directly
with the pinball loss to produce probabilistic individual load forecasts.  By
guiding the network to minimise $\mathbb{E}[\mathcal{L}_\tau]$ rather than MSE,
their approach yields calibrated quantile predictions without requiring post-hoc
quantile regression on the residuals.  Their results show improved reliability and
sharpness compared to symmetric-loss baselines.

\textcite{lin2020asymmetric} address the related problem of contract-capacity
optimisation for high-voltage consumers, where the electricity pricing structure
is inherently asymmetric: underestimating demand triggers penalty charges at a
higher per-unit rate than overestimating.  They propose several asymmetric loss
functions derived directly from the Taiwan Power Company's pricing formula, and
demonstrate that LSTM models trained with these loss functions reduce electricity
costs by 0.88\%--2.42\% over symmetric MSE baselines.

\textcite{wu2021svr} extend this idea to support vector regression (SVR),
developing a cost-oriented asymmetric SVR (AsySVR) framework based on a
linear-linear loss function with an insensitivity parameter to improve
generalisation.  Evaluated on New South Wales load data, their asymmetric
framework achieves daily economic cost reductions of 42\%--57\% relative to
standard SVR, underscoring that the choice of loss function can have a large
practical impact on operational outcomes.

\textcite{zhang2022costoriented} further formalise the decision-oriented
perspective on load forecasting, framing the problem explicitly in terms of the
downstream cost incurred by the system operator.  Their cost-oriented framework
bridges the gap between forecast-accuracy metrics and operational objectives, and
provides theoretical justification for replacing standard point-forecast losses
with cost-aware alternatives.

Taken together, this body of work establishes that asymmetric, cost-sensitive
training objectives are both theoretically motivated and empirically superior to
symmetric alternatives in electricity applications.  Two gaps, however, remain.
First, existing studies almost exclusively employ \emph{piecewise-linear}
asymmetric losses---the pinball loss, or the linear-linear penalties derived from
a specific tariff---and thus inherit the newsvendor model's linearity assumption,
even though real procurement costs are frequently non-linear.  Second, and more
practically, these methods are typically benchmarked against statistical
baselines rather than against the forecast an experienced system operator
actually deploys, leaving open whether cost-aware training improves upon a strong,
real-world operational forecast.

Our work is designed to close both gaps.  We formulate the forecasting problem
with a general, not-necessarily-linear asymmetric cost function that subsumes the
pinball loss as a special case, and train a neural network to minimise it
end-to-end.  Crucially, we evaluate against the \emph{real} day-ahead forecasts
published by the Israeli system operator Noga---an operational baseline with
access to proprietary data and domain expertise---rather than against a synthetic
statistical model.  This lets us ask not merely whether cost-aware training beats
a symmetric loss in the abstract, but whether it can improve upon the forecast
currently used to help run a national grid.  The cost-minimising quantile that
this training targets is characterised formally in Section~\ref{sec:modeling}.

\subsection{Contributions}

The contributions of this paper are as follows:

\begin{enumerate}
    \item \textbf{A general asymmetric-cost formulation.}  We formulate day-ahead
          demand forecasting under a general, not-necessarily-linear asymmetric
          cost function that subsumes the pinball loss as a special case, and we
          train a feedforward neural network to minimise it end-to-end
          (Sections~\ref{sec:modeling}--\ref{sec:method}).

    \item \textbf{Evaluation against a real operational forecast.}  Rather than
          benchmarking against a synthetic statistical baseline, we compare
          against the \emph{actual} day-ahead forecasts published by the Israeli
          system operator Noga --- a strong, real-world baseline produced with
          proprietary data and domain expertise.

    \item \textbf{A practical model that improves on the operator's forecast.}
          Although raw predictive accuracy was \emph{not} a goal of this work,
          our model outperforms the Noga forecast not only on cost-sensitive
          metrics but even under symmetric error metrics such as mean squared
          error.  We regard this as an incidental by-product of a deliberately
          simple architecture and believe that substantial further improvement
          is readily achievable.

    \item \textbf{A demonstration of economic impact.}  We quantify the
          operational cost reduction that cost-aware training delivers over
          symmetric baselines across a range of cost scenarios, showing an effect
          that is large enough --- and cheap enough to obtain --- to be practical
          to adopt in a live market setting.

    \item \textbf{An open-source implementation.}  We release our complete code
          and data pipeline, from raw data ingestion through model training and
          evaluation, as open source,\footnote{\url{https://github.com/gkutiel/noga}}
          so that our results can be reproduced and extended.
\end{enumerate}
